\documentclass[../main.tex]{subfiles}

\begin{document}

\chapter{微积分}\label{chapter:appC-微积分}

\section{边界}\label{section:appC-边界}

\begin{definition}\label{definition:appC-Ck边界}
设 $U\subset\mathbb{R}^n$ 开且有界，$k\in\{1,2,\cdots\}$。称边界 $\partial U$ 是 $C^k$ 的，若对每个点 $x_0\in\partial U$,存在 $r>0$，正交矩阵 $Q$，以及 $C^k$ 函数 $\gamma:\mathbb{R}^{n-1}\to\mathbb{R}$，使得在坐标
\begin{align*}
y\triangleq Q(x-x_0)
\end{align*}
下，有
\begin{align*}
U\cap B(0,r)=\{y\in B(0,r)\mid y_n>\gamma(y_1,\cdots,y_{n-1})\}.
\end{align*}
类似地，$\partial U$ 是 $C^\infty$ 的，若 $\partial U$ 对 $k=1,2,\cdots$ 都是 $C^k$ 的；$\partial U$ 是解析的，若映射 $\gamma$ 解析。
\end{definition}
\begin{remark}
这里 $Q$ 可取为最后一行为 $x_0$ 处关于 $U$ 的单位外法向量$\nu(x_0)$的转置$\nu(x_)^T$正交矩阵.
\end{remark}

\begin{definition}\label{definition:appC-法向量}
\begin{enumerate}[(i)]
\item 若 $\partial U$ 是 $C^1$ 的，则沿 $\partial U$ 定义了单位外法向量场
\begin{align*}
\boldsymbol{\nu}=(\nu_1,\cdots,\nu_n).
\end{align*}
任一点 $x_0\in\partial U$ 处的单位法向量为 $\boldsymbol{\nu}(x_0)=\nu=(\nu_1,\cdots,\nu_n)$。

\item 设 $u\in C^1(\overline{U})$，$r$是单位向量，则称$r\cdot Du\triangleq \frac{\partial u}{\partial r}$为$u$在方向$r$上的方向导，特别地，我们记
\begin{align*}
\frac{\partial u}{\partial\nu}\coloneqq\boldsymbol{\nu}\cdot Du
\end{align*}
为 $u$ 的（外）法向导数。
\end{enumerate}
\end{definition}

\textbf{拉直边界。} 我们经常需要在 $\partial U$ 某点附近改变坐标以"拉平"边界。更具体地，固定 $x_0\in\partial U$，如上选取 $r,\gamma$ 等。定义
\begin{align*}
\begin{cases}
y_i=x_i\coloneqq\Phi_i(x)\quad (i=1,\cdots,n-1),\\
y_n=x_n-\gamma(x_1,\cdots,x_{n-1})\coloneqq\Phi_n(x),
\end{cases}
\end{align*}
写 $y=\boldsymbol{\Phi}(x)$。类似地设
\begin{align*}
\begin{cases}
x_i=y_i\coloneqq\Psi_i(y)\quad (i=1,\cdots,n-1),\\
x_n=y_n+\gamma(y_1,\cdots,y_{n-1})\coloneqq\Psi_n(y),
\end{cases}
\end{align*}
写 $x=\boldsymbol{\Psi}(y)$。则 $\boldsymbol{\Phi}=\boldsymbol{\Psi}^{-1}$，映射 $x\mapsto\boldsymbol{\Phi}(x)=y$ 在 $x_0$ 附近"拉直 $\partial U$"。还注意 $\det D\boldsymbol{\Phi}=\det D\boldsymbol{\Psi}=1$。

\section{Gauss--Green 定理}\label{section:appC-GaussGreen}

本节假设 $U$ 是 $\mathbb{R}^n$ 的有界开子集，$\partial U$ 是 $C^1$ 的。

\begin{theorem}[Gauss--Green 定理。]\label{theorem:appC-GaussGreen}
\begin{enumerate}[(i)]
\item 设 $u\in C^1(\overline{U})$。则
\begin{align}
\int_U u_{x_i}\,\mathrm{d}x=\int_{\partial U}u\nu_i\,\mathrm{d}S \quad (i=1,\cdots,n).
\label{align:appC-GG1}
\end{align}
\item 有
\begin{align}
\int_U\operatorname{div}\mathbf{u}\,\mathrm{d}x=\int_{\partial U}\mathbf{u}\cdot\boldsymbol{\nu}\,\mathrm{d}S
\label{align:appC-GG2}
\end{align}
对每个向量场 $\mathbf{u}\in C^1(\overline{U};\mathbb{R}^n)$。
\end{enumerate}
\end{theorem}

断言 (ii) 也称为散度定理，由 (i) 应用于 $\mathbf{u}=(u_1,\cdots,u_n)$ 的每个分量得到。

\begin{theorem}[分部积分公式。]\label{theorem:appC-分部积分}设 $u,v\in C^1(\overline{U})$。则
\begin{align}
\int_U u_{x_i}v\,\mathrm{d}x=-\int_U uv_{x_i}\,\mathrm{d}x+\int_{\partial U}uv\nu_i\,\mathrm{d}S \quad (i=1,\cdots,n).
\label{align:appC-分部积分}
\end{align}
\end{theorem}

\begin{proof}
对 $uv$ 应用\cref{theorem:appC-GaussGreen}(i)。
\end{proof}

\begin{theorem}[Green 公式。]\label{theorem:appC-Green公式}设 $u,v\in C^2(\overline{U})$。则
\begin{enumerate}[(i)]
\item $\displaystyle\int_U\Delta u\,\mathrm{d}x=\int_{\partial U}\frac{\partial u}{\partial\nu}\,\mathrm{d}S$。
\item $\displaystyle\int_U Dv\cdot Du\,\mathrm{d}x=-\int_U u\Delta v\,\mathrm{d}x+\int_{\partial U}\frac{\partial v}{\partial\nu}u\,\mathrm{d}S$。
\item $\displaystyle\int_U u\Delta v-v\Delta u\,\mathrm{d}x=\int_{\partial U}u\frac{\partial v}{\partial\nu}-v\frac{\partial u}{\partial\nu}\,\mathrm{d}S$。
\end{enumerate}
\end{theorem}

\begin{proof}
利用\eqref{align:appC-分部积分}，以 $u_{x_i}$ 代替 $u$，$v\equiv 1$，得
\begin{align*}
\int_U u_{x_ix_i}\,\mathrm{d}x=\int_{\partial U}u_{x_i}\nu_i\,\mathrm{d}S.
\end{align*}
对 $i=1,\cdots,n$ 求和建立 (i)。

为推导 (ii)，以 $v_{x_i}$ 代替 $v$ 应用\eqref{align:appC-分部积分}。交换 $u$ 和 $v$ 写 (ii) 然后相减，得到 (iii)。
\end{proof}

\section{极坐标、余面积公式}\label{section:appC-极坐标}

接下来将 $n$ 维积分转化为球面上的积分。

\begin{theorem}[极坐标公式]\label{theorem:appC-极坐标}
\begin{enumerate}[(i)]
\item 设 $f:\mathbb{R}^n\to\mathbb{R}$ 连续且可积。则
\begin{align*}
\int_{\mathbb{R}^n}f\,\mathrm{d}x=\int_0^\infty\left(\int_{\partial B(x_0,r)}f\,\mathrm{d}S\right)\mathrm{d}r
\end{align*}
对每个点 $x_0\in\mathbb{R}^n$。
\item 特别地
\begin{align*}
\frac{\mathrm{d}}{\mathrm{d}r}\left(\int_{B(x_0,r)}f\,\mathrm{d}x\right)=\int_{\partial B(x_0,r)}f\,\mathrm{d}S
\end{align*}
对每个 $r>0$。
\end{enumerate}
\end{theorem}

\cref{theorem:appC-极坐标} 是以下定理的特殊情形。

\begin{theorem}[余面积公式。]\label{theorem:appC-余面积}设 $u:\mathbb{R}^n\to\mathbb{R}$ Lipschitz 连续，假设对 a.e. $r\in\mathbb{R}$，水平集
\begin{align*}
\{x\in\mathbb{R}^n\mid u(x)=r\}
\end{align*}
是 $\mathbb{R}^n$ 中光滑 $(n-1)$ 维超曲面。又设 $f:\mathbb{R}^n\to\mathbb{R}$ 连续且可积。则
\begin{align*}
\int_{\mathbb{R}^n}f|Du|\,\mathrm{d}x=\int_{-\infty}^\infty\left(\int_{\{u=r\}}f\,\mathrm{d}S\right)\mathrm{d}r.
\end{align*}
\end{theorem}

\cref{theorem:appC-极坐标} 由\cref{theorem:appC-余面积} 取 $u(x)=|x-x_0|$ 得到。更多余面积公式内容见 \cite{E-G} 第 3 章。

\section{移动区域}\label{section:appC-移动区域}

考虑一族光滑有界区域 $U(\tau)\subset\mathbb{R}^n$ 光滑依赖于参数 $\tau\in\mathbb{R}$。写 $\mathbf{v}$ 为移动边界 $\partial U(\tau)$ 的速度，$\boldsymbol{\nu}$ 为外指单位法向量。

\begin{theorem}[移动区域的微分公式。]\label{theorem:appC-移动区域微分}若 $f=f(x,\tau)$ 是光滑函数，则
\begin{align*}
\frac{\mathrm{d}}{\mathrm{d}\tau}\int_{U(\tau)}f\,\mathrm{d}x=\int_{\partial U(\tau)}f\,\mathbf{v}\cdot\boldsymbol{\nu}\,\mathrm{d}S+\int_{U(\tau)}f_\tau\,\mathrm{d}x.
\end{align*}
\end{theorem}

证明见 Chorin--Marsden~\cite{C-M}。

\section{卷积与磨光}\label{section:appC-卷积磨光}

\begin{definition}\label{notation:appC-Ueps}
设$U\subset\mathbb{R}^n$ 开且 $\varepsilon>0$，记
\begin{align*}
U_\varepsilon\coloneqq\{x\in U\mid\operatorname{dist}(x,\partial U)>\varepsilon\}.
\end{align*}
\end{definition}

\begin{definition}\label{definition:appC-磨光子}
\begin{enumerate}[(i)]
\item 定义 $\eta\in C^\infty(\mathbb{R}^n)$ 为
\begin{align*}
\eta(x)\coloneqq
\begin{cases}
C\exp\left(\dfrac{1}{|x|^2-1}\right) & \text{若 } |x|<1,\\
0 & \text{若 } |x|\geqslant 1,
\end{cases}
\end{align*}
常数 $C>0$ 选取得使 $\int_{\mathbb{R}^n}\eta\,\mathrm{d}x=1$。称 $\eta$ 为\textbf{标准光滑化算子}或\textbf{标准磨光子}。

\item 对每个 $\varepsilon>0$，设
\begin{align*}
\eta_\varepsilon(x)\coloneqq\frac{1}{\varepsilon^n}\eta\left(\frac{x}{\varepsilon}\right).
\end{align*}
函数 $\eta_\varepsilon$ 是 $C^\infty$ 的，满足
\begin{align*}
\int_{\mathbb{R}^n}\eta_\varepsilon\,\mathrm{d}x=1,\quad\operatorname{spt}(\eta_\varepsilon)\subset B(0,\varepsilon).
\end{align*}
\end{enumerate}
\end{definition}

\begin{definition}\label{definition:appC-磨光}
设$U\subset\mathbb{R}^n$ 开且 $\varepsilon>0$，若 $f:U\to\mathbb{R}$ 局部可积，定义其光滑化
\begin{align*}
f_\varepsilon\coloneqq\eta_\varepsilon*f,\quad \text{在} U_\varepsilon \text{中}.
\end{align*}
即
\begin{align*}
f_\varepsilon(x)=\int_U\eta_\varepsilon(x-y)f(y)\,\mathrm{d}y=\int_{B(0,\varepsilon)}\eta_\varepsilon(y)f(x-y)\,\mathrm{d}y
\end{align*}
对 $x\in U_\varepsilon$。
\end{definition}

\begin{theorem}[磨光子的性质]\label{theorem:appC-光滑算子的性质}
设 $U\subseteq\mathbb{R}^n$ 为开集, $f\in L^1_{\mathrm{loc}}(U)$,$\eta\in C_c^\infty(\mathbb{R}^n)$是标准磨光子满足
\begin{align*}
\eta\geqslant0,\quad
\operatorname{spt}\eta\subseteq\overline{B(0,1)},\quad
\int_{\mathbb{R}^n}\eta(x)\,\mathrm{d}x=1.
\end{align*}
对于每个 $\varepsilon>0$, 定义
\begin{align*}
\eta_\varepsilon(x)
&\triangleq\frac{1}{\varepsilon^n}\eta\left(\frac{x}{\varepsilon}\right),\quad
U_\varepsilon
\triangleq\{x\in U\mid\operatorname{dist}(x,\partial U)>\varepsilon\},\\
f_\varepsilon(x)
&\triangleq
\begin{cases}
(\eta_\varepsilon*f)(x)
=\displaystyle\int_U\eta_\varepsilon(x-y)f(y)\,\mathrm{d}y,
& x\in U_\varepsilon,\\
0, & x\in\mathbb{R}^n\setminus U_\varepsilon.
\end{cases}
\end{align*}
则有以下结论, 其中所有收敛均指 $\varepsilon\to0^+$.
\begin{enumerate}[(1)]
\item\label{theorem:appC-光滑算子的性质-1}
$f_\varepsilon\in C^\infty(U_\varepsilon)$, 且对每个多重指标 $\alpha$, 有
\begin{align*}
D^{\alpha}f_{\varepsilon}(x)=\int_U{(D^{\alpha}\eta _{\varepsilon})(x-y)f(y)\,\mathrm{d}y},\quad x\in U_{\varepsilon}.
\end{align*}

\item\label{theorem:appC-光滑算子的性质-2} $f_\varepsilon(x)\to f(x)$ 对几乎处处的 $x\in U$ 成立.

\item\label{theorem:appC-光滑算子的性质-3} 若 $f\in C(U)$, 则 $f_\varepsilon\to f$ 在 $U$ 的任意紧子集上一致成立.

\item\label{theorem:appC-光滑算子的性质-4} 若 $1\leqslant p<\infty$ 且 $f\in L^p_{\mathrm{loc}}(U)$, 则 $f_\varepsilon\to f$ 于 $L^p_{\mathrm{loc}}(U)$.

\item\label{theorem:appC-光滑算子的性质-6}
记 $\operatorname{spt}f$ 为 $f$的支集, 则
\begin{align*}
\operatorname{spt}f_\varepsilon
\subseteq\operatorname{spt}f+\operatorname{spt}\eta_\varepsilon\subseteq\operatorname{spt}f+\overline{B(0,\varepsilon)}=\{x\in\mathbb{R}^n\mid
\operatorname{dist}(x,\operatorname{spt}f)\leqslant\varepsilon\}.
\end{align*}
特别地, 若 $f$ 在 $U$ 内具有紧支集, 且开集 $W$ 满足
\begin{align*}
\operatorname{spt}f\subset\subset W\subset\subset U,
\end{align*}
则当 $\varepsilon>0$ 充分小时, 有
\begin{align*}
\operatorname{spt}f_\varepsilon\subseteq W,\quad
f_\varepsilon\in C_c^\infty(W).
\end{align*}

\item\label{theorem:appC-光滑算子的性质-5} 若 $1\leqslant p<\infty$ 且 $f\in L^p(U)$, 则$\|f_\varepsilon-f\|_{L^p(U)}\to0.$
\end{enumerate}
\end{theorem}
\begin{proof}

由于 \(\eta\in C_c^\infty(\mathbb R^n)\)，\(\operatorname{spt}\eta\) 为紧集。对任意多重指标 \(\alpha\)，有
\[
D^\alpha\eta\in C_c^\infty(\mathbb R^n),
\qquad
\operatorname{spt}(D^\alpha\eta)\subseteq \operatorname{spt}\eta.
\]
因此 \(D^\alpha\eta\) 连续且具有紧支集。连续函数在紧集上有界，故
\[
\|D^\alpha\eta\|_{L^\infty(\mathbb R^n)}
=\sup_{x\in \operatorname{spt}(D^\alpha\eta)}|D^\alpha\eta(x)|
<+\infty.
\]
特别地，取 \(\alpha=0\)，得到
\[
\|\eta\|_{L^\infty(\mathbb R^n)}<+\infty.
\]
又由
\[
\eta_\varepsilon(x)=\frac1{\varepsilon^n}\eta\left(\frac{x}{\varepsilon}\right)
\]
可知，对任意多重指标 \(\alpha\)，有
\[
D^\alpha\eta_\varepsilon(x)
=\varepsilon^{-n-|\alpha|}(D^\alpha\eta)\left(\frac{x}{\varepsilon}\right),
\qquad x\in\mathbb R^n.
\]
于是
\[
\|D^\alpha\eta_\varepsilon\|_{L^\infty(\mathbb R^n)}
=\varepsilon^{-n-|\alpha|}\|D^\alpha\eta\|_{L^\infty(\mathbb R^n)}
<+\infty.
\]

\textbf{(1)}
固定 \(x_0\in U_\varepsilon\)。由 \(U_\varepsilon\) 的定义，存在 \(r>0\)，使得
\[
\varepsilon<r<\operatorname{dist}(x_0,\partial U),
\]
从而 \(\overline{B(x_0,r)}\subset U\)。取开集 \(V\) 使得
\[
\overline{B(x_0,r)}\subset V\subset\subset U.
\]
当 \(|h|\) 充分小时，\(x_0+h e_i\in U_\varepsilon\)，并且
\[
B(x_0,\varepsilon)\cup B(x_0+h e_i,\varepsilon)\subset B(x_0,r)\subset V.
\]
于是
\[
\frac{f_\varepsilon(x_0+h e_i)-f_\varepsilon(x_0)}{h}
=\int_U \frac{\eta_\varepsilon(x_0+h e_i-y)-\eta_\varepsilon(x_0-y)}{h}f(y)\,\mathrm dy.
\]
由于上述被积函数的支集含于 \(V\)，可将积分区域改为 \(V\)。又由微积分基本定理知
\[
\frac{\eta_\varepsilon(x_0+h e_i-y)-\eta_\varepsilon(x_0-y)}{h}
=\int_0^1 \partial_i\eta_\varepsilon(x_0+t h e_i-y)\,\mathrm dt,
\]
且当 \(h\to0\) 时，有
\[
\partial_i\eta_\varepsilon(x_0+t h e_i-y)\to \partial_i\eta_\varepsilon(x_0-y)
\]
关于 \(y\in V\) 一致成立，同时有
\[
\left|\int_0^1 \partial_i\eta_\varepsilon(x_0+t h e_i-y)\,\mathrm dt\right|
\leqslant \|\partial_i\eta_\varepsilon\|_{L^{\infty}\left( \mathbb{R} ^n \right) },
\]
而 \(f\in L^1_{\mathrm{loc}}(U)\)，故 \(f\in L^1(V)\)。由\hyperref[Real Analysis-theorem:控制收敛定理]{控制收敛定理}，有
\[
\partial_i f_\varepsilon(x_0)
=\int_U \partial_i\eta_\varepsilon(x_0-y)f(y)\,\mathrm dy.
\]
因此 \(f_\varepsilon\) 在 \(U_\varepsilon\) 内一阶偏导数存在，且公式成立。对多重指标 \(\alpha\) 反复应用同样论证，即得
\[
D^\alpha f_\varepsilon(x)
=\int_U D^\alpha\eta_\varepsilon(x-y)f(y)\,\mathrm dy,
\qquad x\in U_\varepsilon.
\]
右端为含参变量积分，被积函数关于 \(x\) 连续，故 \(D^\alpha f_\varepsilon\in C(U_\varepsilon)\)。因此
\[
f_\varepsilon\in C^\infty(U_\varepsilon).
\]

\textbf{(2)}
由\hyperref[Real Analysis-corollary:可积函数的Lebesgue点]{Lebesgue微分定理}，对几乎处处的 \(x\in U\)，有
\[
\lim_{r\to0}\fint_{B(x,r)}|f(y)-f(x)|\,\mathrm dy=0.
\]
固定这样的 \(x\)。当 \(\varepsilon<\operatorname{dist}(x,\partial U)\) 时，\(x\in U_\varepsilon\)，且 \(B(x,\varepsilon)\subset U\)。由 \(\int_{\mathbb R^n}\eta=1\)，有
\[
\int_U \eta_\varepsilon(x-y)\,\mathrm dy=1.
\]
因此
\[
f_\varepsilon(x)-f(x)
=\int_U \eta_\varepsilon(x-y)\bigl(f(y)-f(x)\bigr)\,\mathrm dy.
\]
于是
\[
\begin{aligned}
|f_\varepsilon(x)-f(x)|
&\leqslant \int_{B(x,\varepsilon)}\eta_\varepsilon(x-y)|f(y)-f(x)|\,\mathrm dy\leqslant \frac{\|\eta\|_{L^{\infty}\left( \mathbb{R} ^n \right) }}{\varepsilon^n}
\int_{B(x,\varepsilon)}|f(y)-f(x)|\,\mathrm dy\\
&=\|\eta\|_{L^{\infty}\left( \mathbb{R} ^n \right) }|B(0,1)|
\fint_{B(x,\varepsilon)}|f(y)-f(x)|\,\mathrm dy\to0,
\qquad \varepsilon\to0^+.
\end{aligned}
\]
故 \(f_\varepsilon(x)\to f(x)\) 对几乎处处的 \(x\in U\) 成立。

\textbf{(3)}
设 \(f\in C(U)\)，并设 \(K\subset\subset U\) 为紧集。取开集 \(W\) 使得
\[
K\subset W\subset\subset U.
\]
则 \(f\) 在紧集 \(\overline W\) 上一致连续。任给 \(\delta>0\)，存在 \(\rho>0\)，使得当 \(x,y\in\overline W\) 且 \(|x-y|<\rho\) 时，
\[
|f(x)-f(y)|<\delta.
\]
取 \(\varepsilon>0\) 充分小，使得
\[
\varepsilon<\rho,\qquad \varepsilon<\operatorname{dist}(K,\partial W),\qquad \varepsilon<\operatorname{dist}(K,\partial U).
\]
则对每个 \(x\in K\)，有 \(B(x,\varepsilon)\subset W\)，且 \(x\in U_\varepsilon\)。由 (2) 中的估计，
\[
|f_\varepsilon(x)-f(x)|
\leqslant \|\eta\|_{L^\infty}|B(0,1)|
\fint_{B(x,\varepsilon)}|f(y)-f(x)|\,\mathrm dy
\leqslant C\delta.
\]
其中 \(C=\|\eta\|_{L^\infty}|B(0,1)|\)。由于 \(\delta>0\) 任意，故
\[
f_\varepsilon\to f
\]
在 \(K\) 上一致成立。由于 \(K\subset\subset U\) 任意，故 \(f_\varepsilon\to f\) 在 \(U\) 的任意紧子集上一致成立。

\textbf{(4)}
设 \(1\leqslant p<\infty\)，\(f\in L^p_{\mathrm{loc}}(U)\)。先证明局部 \(L^p\) 估计。取开集 \(V,W\) 使得
\[
V\subset\subset W\subset\subset U.
\]
当 \(\varepsilon<\operatorname{dist}(V,\partial W)\) 时，对 \(x\in V\)，有 \(B(x,\varepsilon)\subset W\)。由\hyperref[Real Analysis-theorem:实变函数--Hölder不等式]{Hölder 不等式}，得
\[
\begin{aligned}
|f_\varepsilon(x)|
&\leqslant \int_{B(x,\varepsilon)}\eta_\varepsilon(x-y)|f(y)|\,\mathrm dy\\
&\leqslant \left(\int_{B(x,\varepsilon)}\eta_\varepsilon(x-y)\,\mathrm dy\right)^{1-\frac{1}{p}}
\left(\int_{B(x,\varepsilon)}\eta_\varepsilon(x-y)|f(y)|^p\,\mathrm dy\right)^{\frac{1}{p}}\\
&=\left(\int_{B(x,\varepsilon)}\eta_\varepsilon(x-y)|f(y)|^p\,\mathrm dy\right)^{\frac{1}{p}}.
\end{aligned}
\]
故
\[
|f_\varepsilon(x)|^p
\leqslant \int_{B(x,\varepsilon)}\eta_\varepsilon(x-y)|f(y)|^p\,\mathrm dy.
\]
对 \(x\in V\) 积分，并利用 \hyperref[Real Analysis-theorem:实变函数-Fubini定理]{Fubini定理}，得
\[
\begin{aligned}
\int_V |f_\varepsilon(x)|^p\,\mathrm dx
&\leqslant \int_V\int_W \eta_\varepsilon(x-y)|f(y)|^p\,\mathrm dy\,\mathrm dx\\
&=\int_W |f(y)|^p\left(\int_V \eta_\varepsilon(x-y)\,\mathrm dx\right)\mathrm dy\\
&\leqslant \int_W |f(y)|^p\,\mathrm dy.
\end{aligned}
\]
因此
\[
\|f_\varepsilon\|_{L^p(V)}\leqslant \|f\|_{L^p(W)}.
\]

现在给定 \(K\subset\subset U\)，取开集 \(V,W\) 使得
\[
K\subset V\subset\subset W\subset\subset U.
\]
任给 \(\delta>0\)。由于 \(C_c(W)\) 在 \(L^p(W)\) 中稠密，存在 \(g\in C_c(W)\)，使得
\[
\|f-g\|_{L^p(W)}<\delta.
\]
将 \(g\) 零延拓到 \(U\)。由 (3)，\(g_\varepsilon\to g\) 在 \(V\) 上一致，故
\[
\|g_\varepsilon-g\|_{L^p(V)}\to0.
\]
于是
\[
\begin{aligned}
\|f_\varepsilon-f\|_{L^p(K)}
&\leqslant \|f_\varepsilon-f\|_{L^p(V)}\leqslant \|f_\varepsilon-g_\varepsilon\|_{L^p(V)}
+\|g_\varepsilon-g\|_{L^p(V)}
+\|g-f\|_{L^p(V)}\\
&\leqslant \|f-g\|_{L^p(W)}
+\|g_\varepsilon-g\|_{L^p(V)}
+\|g-f\|_{L^p(W)}\\
&\leqslant 2\delta+\|g_\varepsilon-g\|_{L^p(V)}.
\end{aligned}
\]
令 \(\varepsilon\to0^+\)，得
\[
\limsup_{\varepsilon\to0^+}\|f_\varepsilon-f\|_{L^p(K)}\leqslant 2\delta.
\]
再令 \(\delta\to0^+\)，即得
\[
\|f_\varepsilon-f\|_{L^p(K)}\to0.
\]
由于 \(K\subset\subset U\) 任意，故 \(f_\varepsilon\to f\) 于 \(L^p_{\mathrm{loc}}(U)\)。

\textbf{(5)}
记 $S_\eta=\operatorname{spt}\eta$ , $S_f=\operatorname{spt}f$ .
由于 $\eta\in C_c(\mathbb R^n)$ , 所以 $S_\eta$ 是紧集 ; 而 $S_f$ 按定义是闭集 .

首先证明 $S_\eta+S_f$ 是闭集 . 设
\begin{align*}
x_k=a_k+b_k,\qquad a_k\in S_\eta,\quad b_k\in S_f,
\end{align*}
且 $x_k\to x$ . 因为 $S_\eta$ 紧 , 存在子列 $a_{k_j}\to a\in S_\eta$ . 于是
\begin{align*}
b_{k_j}=x_{k_j}-a_{k_j}\to x-a.
\end{align*}
又因为 $S_f$ 闭 , 所以 $x-a\in S_f$ . 因此
\begin{align*}
x=a+(x-a)\in S_\eta+S_f.
\end{align*}
故 $S_\eta+S_f$ 是闭集 .

现在任取
\begin{align*}
x\notin S_\eta+S_f.
\end{align*}
对任意 $y\in S_f$ , 若 $x-y\in S_\eta$ , 则
\begin{align*}
x=(x-y)+y\in S_\eta+S_f,
\end{align*}
与 $x\notin S_\eta+S_f$ 矛盾 . 因此必有
\begin{align*}
x-y\notin S_\eta.
\end{align*}
由于 $S_\eta=\operatorname{spt}\eta$ , 故
\begin{align*}
\eta(x-y)=0.
\end{align*}
另一方面 , 若 $y\notin S_f=\operatorname{spt}f$ , 则按支集定义有
\begin{align*}
f(y)=0.
\end{align*}
于是对任意 $y\in\mathbb R^n$ , 都有
\begin{align*}
\eta(x-y)f(y)=0.
\end{align*}
因此
\begin{align*}
(\eta*f)(x)=\int_{\mathbb R^n}\eta(x-y)f(y)\mathrm{d}y=0.
\end{align*}

由于 $S_\eta+S_f$ 是闭集 , 且 $x\notin S_\eta+S_f$ , 存在 $\delta>0$ , 使得
\begin{align*}
B(x,\delta)\cap (S_\eta+S_f)=\varnothing.
\end{align*}
对任意 $z\in B(x,\delta)$ , 同样有 $z\notin S_\eta+S_f$ . 由上面的论证可得
\begin{align*}
(\eta*f)(z)=0.
\end{align*}
所以 $x$ 的某个邻域内 $\eta*f$ 恒为零 , 故
\begin{align*}
x\notin \operatorname{spt}(\eta*f).
\end{align*}
由 $x\notin S_\eta+S_f$ 的任意性 , 得到
\begin{align*}
\operatorname{spt}(\eta*f)\subseteq S_\eta+S_f=\operatorname{spt}\eta+\operatorname{spt}f.
\end{align*}
由于 \(f_\varepsilon\) 在 \(U_\varepsilon\) 内等于 \(\eta_\varepsilon*f\)，在 \(U_\varepsilon\) 外为 \(0\)，故
\[
\operatorname{spt}f_\varepsilon
\subset \operatorname{spt}f+\operatorname{spt}\eta_\varepsilon.
\]
又因为
\[
\operatorname{spt}\eta_\varepsilon
=\varepsilon\operatorname{spt}\eta
\subset \overline{B(0,\varepsilon)},
\]
所以
\[
\operatorname{spt}f_\varepsilon
\subset \operatorname{spt}f+\overline{B(0,\varepsilon)}
=\{x\in\mathbb R^n\mid \operatorname{dist}(x,\operatorname{spt}f)\le\varepsilon\}.
\]

特别地，若 \(f\) 在 \(U\) 内具有紧支集，且开集 \(W\) 满足
\[
\operatorname{spt}f\subset\subset W\subset\subset U,
\]
则由于 \(\operatorname{spt}f\) 紧且含于 \(W\)，有
\[
\operatorname{dist}(\operatorname{spt}f,\partial W)>0.
\]
取 \(\varepsilon>0\) 充分小，使得
\[
\operatorname{spt}f+\overline{B(0,\varepsilon)}\subset W,
\]
同时由于 \(W\subset\subset U\)，还可取 \(\varepsilon<\operatorname{dist}(W,\partial U)\)，从而 \(W\subset U_\varepsilon\)。于是
\[
\operatorname{spt}f_\varepsilon\subset W,
\]
且由 (1) 知 \(f_\varepsilon\in C^\infty(U_\varepsilon)\)，故 \(f_\varepsilon\in C^\infty(W)\)。又 \(\operatorname{spt}f_\varepsilon\subset W\) 且 \(W\) 紧包含于 \(U\)，所以
\[
f_\varepsilon\in C_c^\infty(W).
\]

\textbf{(6)}固定 \(\sigma>0\)。因为 \(f\in L^p(U)\)，所以由\hyperref[Real Analysis-theorem:积分的绝对连续性]{积分的绝对连续性}知,存在紧集 \(K_1\subset\subset U\)，使得
\begin{align*}
\int_{U\setminus K_1}|f|^p\,\mathrm{d}x<\sigma^p.
\end{align*}
取开集 \(V\) 使得 \(K_1\subset\subset V\subset\subset U\)。令 \(K_2\triangleq\overline{V}\)，则 \(K_2\) 为紧集且 \(K_2\subset\subset U\)。于是
\begin{align*}
\int_{U\setminus K_2}|f|^p\,\mathrm{d}x\leqslant\int_{U\setminus K_1}|f|^p\,\mathrm{d}x<\sigma^p.
\end{align*}
由\hyperref[proposition:整体可积与局部可积的关系]{整体可积与局部可积的关系}知$f\in L_{\text{loc}}^p(U)$,故由\hyperref[theorem:appC-光滑算子的性质]{光滑算子性质\ref{theorem:appC-光滑算子的性质-4}}，有
\begin{align*}
\|f_\varepsilon-f\|_{L^p(K_2)}\to0,\quad \varepsilon\to0^+.
\end{align*}
故存在 \(\varepsilon_0>0\)，使得当 \(0<\varepsilon<\varepsilon_0\) 时，
\begin{align}\label{eqHJFIJOFIOWJIOWIJO-1}
\|f_\varepsilon-f\|_{L^p(K_2)}<\sigma.
\end{align}
现在估计 \(U\setminus K_2\) 上的积分。首先
\begin{align}\label{eqHJFIJOFIOWJIOWIJO-2}
\|f_\varepsilon-f\|_{L^p(U\setminus K_2)}
\leqslant\|f_\varepsilon\|_{L^p(U\setminus K_2)}+\|f\|_{L^p(U\setminus K_2)}<\|f_\varepsilon\|_{L^p(U\setminus K_2)}+\sigma.
\end{align}
只需控制 \(\|f_\varepsilon\|_{L^p(U\setminus K_2)}\)。由于 \(f_\varepsilon\) 在 \(U_\varepsilon\) 外为零，利用\hyperref[Real Analysis-theorem:Jensen不等式]{Jensen不等式}和\hyperref[Real Analysis-theorem:实变函数-Fubini定理]{Fubini 定理}有
\begin{align}
\|f_\varepsilon\|_{L^p(U\setminus K_2)}^p
&=\int_{(U\setminus K_2)\cap U_\varepsilon}|f_\varepsilon(x)|^p\,\mathrm{d}x\nonumber\\
&\leqslant\int_{(U\setminus K_2)\cap U_\varepsilon}\left(\int_U\eta_\varepsilon(x-y)|f(y)|^p\,\mathrm{d}y\right)\mathrm{d}x\nonumber\\
&=\int_U|f(y)|^p\left(\int_{(U\setminus K_2)\cap U_\varepsilon}\eta_\varepsilon(x-y)\,\mathrm{d}x\right)\mathrm{d}y.\label{HFIOUHJEIOWJIOWJIOJWIODIJOIJOeq-1}
\end{align}
将最后的积分分成 \(y\in K_1\) 和 \(y\in U\setminus K_1\) 两部分。

对于 \(y\in U\setminus K_1\)，内层积分不超过 \(1\)，故
\begin{align}\label{HFIOUHJEIOWJIOWJIOJWIODIJOIJOeq-2}
\int_{U\setminus K_1}|f(y)|^p\left(\int_{(U\setminus K_2)\cap U_\varepsilon}\eta_\varepsilon(x-y)\,\mathrm{d}x\right)\mathrm{d}y
\leqslant\int_{U\setminus K_1}|f(y)|^p\,\mathrm{d}y<\sigma^p.
\end{align}
对于 \(y\in K_1\)，由\subcref{Real Analysis-proposition:点集间的距离的性质}{Real Analysis-proposition:点集间的距离的性质-3}知 \(K_1\) 与 \(U\setminus K_2\) 的距离至少为
\begin{align*}
\rho\triangleq\operatorname{dist}(K_1,U\setminus K_2)>0.
\end{align*}
当 \(0<\varepsilon<\rho\)且\(x\in U\setminus K_2\)时，若\(|x-y|<\varepsilon\)，则 \(x\) 与 \(K_1\) 的距离小于 \(\varepsilon<\rho\)，从而 \(x\in K_2\)，这与 \(x\in U\setminus K_2\) 矛盾!因此当 \(0<\varepsilon<\rho\)且\(x\in U\setminus K_2\)时，\(|x-y|\geqslant\varepsilon\),而$\text{spt}\,\eta_\varepsilon\subseteq \overline{B(0,\varepsilon)}$,故此时$\eta_\varepsilon(x-y)=0,\,\text{a.e.}\,x\in U\setminus K_1$.进而
\begin{align}\label{HFIOUHJEIOWJIOWJIOJWIODIJOIJOeq-3}
\int_{(U\setminus K_2)\cap U_\varepsilon}\eta_\varepsilon(x-y)\,\mathrm{d}x=0,\quad y\in K_1.
\end{align}
于是当 \(0<\varepsilon<\min\{\varepsilon_0,\rho\}\) 时，结合\eqref{HFIOUHJEIOWJIOWJIOJWIODIJOIJOeq-1}\eqref{HFIOUHJEIOWJIOWJIOJWIODIJOIJOeq-2}\eqref{HFIOUHJEIOWJIOWJIOJWIODIJOIJOeq-3}式有
\begin{align*}
\|f_\varepsilon\|_{L^p(U\setminus K_2)}^p
\leqslant\int_{U\setminus K_1}|f(y)|^p\,\mathrm{d}y<\sigma^p\Longrightarrow \|f_\varepsilon\|_{L^p(U\setminus K_2)}<\sigma.
\end{align*}
因此再由\eqref{eqHJFIJOFIOWJIOWIJO-2}式知
\begin{align}\label{eqHJFIJOFIOWJIOWIJO-3}
\|f_\varepsilon-f\|_{L^p(U\setminus K_2)}<2\sigma.
\end{align}
综合起来，当 \(0<\varepsilon<\min\{\varepsilon_0,\rho\}\) 时，结合\eqref{eqHJFIJOFIOWJIOWIJO-1}\eqref{eqHJFIJOFIOWJIOWIJO-3}式知
\begin{align*}
\|f_\varepsilon-f\|_{L^p(U)}^p
=\|f_\varepsilon-f\|_{L^p(K_2)}^p+\|f_\varepsilon-f\|_{L^p(U\setminus K_2)}^p<\sigma^p+(2\sigma)^p
\leqslant(1+2^p)\sigma^p.
\end{align*}
由于 \(\sigma>0\) 任意，令 \(\sigma\to0^+\)，即得
\begin{align*}
\|f_\varepsilon-f\|_{L^p(U)}\to0,\quad \varepsilon\to0^+.
\end{align*}
\end{proof}

\section{反函数定理}\label{section:appC-反函数}

设 $U\subset\mathbb{R}^n$ 开集，假设 $\mathbf{f}:U\to\mathbb{R}^n$ 是 $C^1$ 的，$\mathbf{f}=(f_1,\cdots,f_n)$。设 $x_0\in U$，$z_0=\mathbf{f}(x_0)$。回忆 \S A.4 中我们写
\begin{align*}
D\mathbf{f}=\begin{pmatrix}
f^1_{x_1} & \cdots & f^1_{x_n}\\
& \ddots & \\
f^n_{x_1} & \cdots & f^n_{x_n}
\end{pmatrix}_{n\times n}=\mathbf{f}\text{ 的梯度矩阵}.
\end{align*}

\begin{definition}\label{definition:appC-Jacobian}
$J\mathbf{f}=$ $\mathbf{f}$ 的 Jacobi 行列式 $=|\det D\mathbf{f}|=\left|\dfrac{\partial(f_1,\cdots,f_n)}{\partial(x_1,\cdots,x_n)}\right|$。
\end{definition}

\begin{theorem}[反函数定理]\label{theorem:appC-反函数}假设 $\mathbf{f}\in C^1(U;\mathbb{R}^n)$ 且 $J\mathbf{f}(x_0)\neq 0$。则存在开集 $V\subset U$（$x_0\in V$）和开集 $W\subset\mathbb{R}^n$（$z_0\in W$）使得
\begin{enumerate}[(i)]
\item 映射 $\mathbf{f}:V\to W$ 是一一对应且到上的，
\item 逆函数 $\mathbf{f}^{-1}:W\to V$ 是 $C^1$ 的。
\item 若 $\mathbf{f}\in C^k$，则 $\mathbf{f}^{-1}\in C^k$（$k=2,\cdots$）。
\end{enumerate}
\end{theorem}

\section{隐函数定理}\label{section:appC-隐函数}

设 $n,m$ 为正整数。我们将 $\mathbb{R}^{n+m}$ 中典型点写为
\begin{align*}
(x,y)=(x_1,\cdots,x_n,y_1,\cdots,y_m)
\end{align*}
其中 $x\in\mathbb{R}^n$，$y\in\mathbb{R}^m$。

设 $U\subset\mathbb{R}^{n+m}$ 开集，假设 $\mathbf{f}:U\to\mathbb{R}^m$ 是 $C^1$ 的，$\mathbf{f}=(f_1,\cdots,f_m)$。设 $(x_0,y_0)\in U$，$z_0=\mathbf{f}(x_0,y_0)$。

\begin{definition}\label{notation:appC-Df}
\begin{align*}
D\mathbf{f}=\begin{pmatrix}
f^1_{x_1} & \cdots & f^1_{x_n} & f^1_{y_1} & \cdots & f^1_{y_m}\\
\vdots & \ddots & \vdots & \vdots & \ddots & \vdots\\
f^m_{x_1} & \cdots & f^m_{x_n} & f^m_{y_1} & \cdots & f^m_{y_m}
\end{pmatrix}_{m\times(n+m)}=(D_x\mathbf{f},D_y\mathbf{f})=\mathbf{f}\text{ 的梯度矩阵}.
\end{align*}
\end{definition}

\begin{definition}\label{definition:appC-Jyf}
$J_y\mathbf{f}=|\det D_y\mathbf{f}|=\left|\dfrac{\partial(f_1,\cdots,f_m)}{\partial(y_1,\cdots,y_m)}\right|$。
\end{definition}

\begin{theorem}[隐函数定理。]\label{theorem:appC-隐函数}假设 $\mathbf{f}\in C^1(U;\mathbb{R}^m)$ 且 $J_y\mathbf{f}(x_0,y_0)\neq 0$。则存在开集 $V\subset U$（$(x_0,y_0)\in V$），开集 $W\subset\mathbb{R}^n$（$x_0\in W$），和 $C^1$ 映射 $\mathbf{g}:W\to\mathbb{R}^m$ 使得
\begin{enumerate}[(i)]
\item $\mathbf{g}(x_0)=y_0$，
\item $\mathbf{f}(x,\mathbf{g}(x))=z_0$（$x\in W$），
\item 若 $(x,y)\in V$ 且 $\mathbf{f}(x,y)=z_0$，则 $y=\mathbf{g}(x)$。
\item 若 $\mathbf{f}\in C^k$，则 $\mathbf{g}\in C^k$（$k=2,\cdots$）。
\end{enumerate}
函数 $\mathbf{g}$ 在 $x_0$ 附近由方程 $\mathbf{f}(x,y)=z_0$ 隐式定义。
\end{theorem}

\section{一致收敛}\label{section:appC-一致收敛}

我们在此记录 Arzel\`a--Ascoli 紧性判别准则用于一致收敛：

假设 $\{f_k\}_{k=1}^\infty$ 是定义在 $\mathbb{R}^n$ 上的实值函数序列，使得
\begin{align*}
|f_k(x)|\leqslant M \quad (k=1,\cdots,\;x\in\mathbb{R}^n)
\end{align*}
\end{document}